\begin{proof}

\begin{enumerate}[label={$\langle1\rangle$\arabic*.},leftmargin=*]
\item Differentiating the three equations gives the presentations in \eqref{wbII:eq:CN:6.1}.  For
$M_2$, the first map in \eqref{wbII:eq:CN:6.2} sends $(a,b,c)$ to $(ax-by,c)$.  If this is
zero, then $c=0$ and $ax=by$.  Since $x$ and $y$ are relatively prime in the
UFD $R_2$, there is an $h$ with $b=xh$ and $a=yh$; hence its kernel is exactly
the displayed relation.  For $M_3$, the second map sends $(a,b,c)$ to
$(ax-cz,by-cz)$.  If it is zero, relative primeness first gives $x\mid c$,
then $y\mid(c/x)$; writing $c=xyh$ gives
$(a,b,c)=h(yz,xz,xy)$, again exactly the defining relation.  Thus both maps
are injective, proving torsion-freeness and \eqref{wbII:eq:CN:6.3}.

\item At the maximal ideal $(x,y,z)$ each defining relation has all coefficients in
that maximal ideal.  Consequently the quotient of either module by the
maximal ideal has dimension $3$ over $\C$, while its generic rank is $2$;
a locally free module of rank $2$ would need only $2$ generators.  This proves
non-local-freeness at the origin.  Finally, for $t=s^m$ one has
$dt=m s^{m-1}ds=0$ relatively.  The integer $m$ is a unit in $\C$, so the
decomposition \eqref{wbII:eq:CN:6.4} follows.  Its last summand is annihilated by
$s^{m-1}$, while the other two summands are free; after inverting $s$ it
vanishes.  It is therefore exactly the torsion submodule.

\end{enumerate}
\end{proof}
